%! Polynomial Functions and Real Zeros — Unit 2, Lesson 2.3 — Lesson Plan
\documentclass[10pt]{article}
\usepackage{precalculus-article}
\usepackage{precalculus-boxes}
\newcommand{\TallMath}[1]{$\displaystyle #1\rule[-1.4em]{0pt}{3.2em}$}
\newcommand{\CourseName}{Precalculus}
\newcommand{\MeetingLength}{55 minutes}
\newcommand{\UnitNumberName}{Unit 2: Polynomial Functions \quad}
\newcommand{\LessonNumberName}{Lesson 2.3: Polynomial Functions and Real Zeros}

\begin{document}
\begin{center}
    {\Large\bfseries \CourseName \\
    {\small\bfseries \UnitNumberName \LessonNumberName}}
\end{center}

\begin{tcolorbox}[colback=lilac, colframe=plum, boxrule=0.9pt, arc=2mm,
    left=2mm, right=2mm, top=1.5mm, bottom=1.5mm]
    \textbf{Primary Objective:} Students will use the factor/zero equivalence
    to read the real zeros of a polynomial from its factored form and to build
    a factored form from given zeros; determine the multiplicity of each real
    zero and predict whether the graph crosses the $x$-axis or is tangent to
    it there; and find the degree of a polynomial from successive differences
    over equally spaced inputs.
    \hfill\textit{\small (CED Topic 1.5, real-zeros half --- Skills 1.B, 2.B)}
\end{tcolorbox}

\renewcommand{\arraystretch}{1.6}
\begin{skillbox}[Priority Ideas \& Skills]{goldbox}
\begin{tabularx}{\linewidth}{@{}X X@{}}
\textbf{AP Skills} &
\textbf{Key Understandings} \\
\midrule
\textbf{1.B} --- Express functions, equations, or expressions in analytically
equivalent forms that are useful in a given context. &
\begin{itemize}[leftmargin=1.2em,itemsep=2pt,topsep=0pt]
  \item If $p(a) = 0$ then $a$ is a zero of $p$ and a root of $p(x) = 0$.
        (1.5.A.1)
  \item For a \textit{real} number $a$: $(x - a)$ is a linear factor of $p$
        \textbf{if and only if} $a$ is a zero of $p$. Both directions matter
        --- read zeros off factors, build factors from zeros. (1.5.A.1)
\end{itemize} \\
\textbf{2.B} --- Construct equivalent graphical, numerical, analytical, and
verbal representations of functions. &
\begin{itemize}[leftmargin=1.2em,itemsep=2pt,topsep=0pt]
  \item A repeated linear factor $(x-a)^n$ gives the zero $a$
        \textit{multiplicity} $n$. (1.5.A.2, real half)
  \item If $a$ is a real zero, the graph has an $x$-intercept at $(a, 0)$;
        real zeros will later serve as endpoints for polynomial inequalities.
        (1.5.A.3)
  \item Even multiplicity: outputs near $x = a$ keep the same sign and the
        graph is \textit{tangent} to the axis. Odd multiplicity: sign changes,
        the graph \textit{crosses}. (1.5.A.5)
  \item Degree from data: over equally spaced inputs, the degree is the least
        $n$ with constant $n$th differences. (1.5.A.6)
\end{itemize} \\
\end{tabularx}
\end{skillbox}

\begin{skillbox}[{Vocabulary, Concepts, \& Theorems}]{greenbox}
\renewcommand{\arraystretch}{1.4}
\begin{tabularx}{\linewidth}{@{} >{\leavevmode\bfseries\color{plum}}p{0.27\linewidth} X @{}}
Zero of a function & A number $a$ with $p(a) = 0$ --- an input whose output is zero. \\
Root               & A solution of the equation $p(x) = 0$; the roots of $p(x)=0$ are exactly the zeros of $p$. \\
$x$-intercept      & A point $(a, 0)$ where the graph meets the $x$-axis; $a$ is a real zero of $p$. \\
Linear factor      & A factor of the form $(x - a)$. For real $a$, it divides $p(x)$ exactly when $a$ is a zero of $p$. \\
Multiplicity       & The number of times the linear factor $(x - a)$ is repeated in the factored form of $p(x)$. \\
Tangent to the $x$-axis & The graph touches the axis at an intercept and turns back without crossing --- the signature of even multiplicity. \\
Successive differences & Differences of differences over equally spaced inputs; the degree is the least $n$ giving a constant $n$th row. \\
\end{tabularx}
\end{skillbox}

\begin{skillbox}[Activate Prior Knowledge \& Spiral Review]{lilac}
    Please see attached warm-up (spiral review).
    Skills reviewed: factoring a quadratic trinomial, solving a factored
    equation with the zero-product property, and evaluating a polynomial at
    a value with the parentheses habit. Item 3's polynomial
    $p(x) = x^3 - 2x^2 - 5x + 6$ returns verbatim in Guided Notes Example 1 ---
    students verify $p(1) = 0$ on the warm-up and $p(3) = 0$ ten minutes
    later, which is the whole zero test.
\end{skillbox}

\begin{skillbox}[Hook]{lilac}
    \textbf{Four questions, one answer.} Put $p(x) = x^2 - x - 6$ on the board
    with four prompts pulled from four different tests:
    \begin{enumerate}[itemsep=2pt]
        \item \textit{Solve $p(x) = 0$.} \quad \textit{Find the roots.} \quad
              \textit{Find the zeros of $p$.} \quad
              \textit{Find the $x$-intercepts.}
        \item Ask groups to decide, without computing, whether these are four
              problems or one problem in four costumes.
    \end{enumerate}
    \textit{Discussion prompt:} they are one problem. The numbers $3$ and
    $-2$ answer all four; only the packaging changes ($x = 3$ as a root,
    $(3, 0)$ as an intercept). Today we get the vocabulary straight and then
    learn what the graph does at each of those numbers.
\end{skillbox}

\begin{skillbox}[Lesson]{lilac}
\begin{multicols}{2}

\textbf{Part 1: Four Names for One Idea} (8 min)

\begin{itemize}
  \item Define: $a$ is a zero of $p$ when $p(a) = 0$. Then walk the table of
        synonyms --- zero (function language), root (equation language),
        $x$-intercept (graph language).
  \item Insist on the format difference: a zero is a \textit{number}, an
        intercept is a \textit{point}. That distinction is worth a point on
        every future test.
  \item Example 1: verify $p(3) = 0$ for $p(x) = x^3 - 2x^2 - 5x + 6$. Then
        say the three conclusions aloud in the three languages.
\end{itemize}

\textbf{Part 2: The Factor/Zero Equivalence} (16 min --- the core)

\begin{itemize}
  \item State it as a two-way street:
        $(x - a)$ is a factor $\Longleftrightarrow$ $a$ is a zero.
  \item Left to right: $p(x) = (x-2)(x+3)(x-5)$ has zeros $2, -3, 5$. Drill
        the sign flip on $(x+3)$ --- this is the most common error of the day.
  \item Right to left: zeros $-1, 0, 4$ of degree 3 give
        $p(x) = x(x+1)(x-4)$. Note that a zero of $0$ contributes the bare
        factor $x$.
  \item Example 2: zeros $2$ and $-4$ with $p(0) = -16$ pins down the leading
        coefficient $a = 2$. This is the AP Skill 1.B move: the factored form
        is chosen because it makes the zeros visible.
\end{itemize}

\columnbreak

\textbf{Part 3: Multiplicity and the Graph} (16 min)

\begin{itemize}
  \item Define multiplicity by counting repeats of a factor; note that the
        multiplicities of the real zeros sum to the degree in every example
        today.
  \item Fill the table for $(x+1)^2(x-3)(x-4)^3$: multiplicity, even/odd,
        crosses/tangent.
  \item Motivate the rule with signs, not memorization: near $x = a$ the
        factor $(x-a)^n$ controls the sign. An even power is never negative,
        so the outputs on both sides match --- the graph must turn back.
  \item Read the pre-drawn graph of $y = (x+1)^2(x-2)$: tangent at $-1$,
        crossing at $2$. Nobody draws anything.
\end{itemize}

\textbf{Part 4: Degree from Successive Differences} (10 min)

\begin{itemize}
  \item Recall 2.1: linear $\Rightarrow$ constant first differences;
        quadratic $\Rightarrow$ constant second. Extend the ladder.
  \item Work the table $5, 4, 9, 26, 61, 120$ row by row on the board:
        $-1, 5, 17, 35, 59$; then $6, 12, 18, 24$; then $6, 6, 6$. Degree 3.
  \item Name the fine print out loud: the inputs must be equally spaced, and
        we need the \textit{least} $n$ with a constant row.
\end{itemize}
\end{multicols}
\end{skillbox}

\begin{skillbox}[Explicit Instruction --- Building a Polynomial from Zeros]{lilac}
Model the four-step routine every time, and require it on the activity and
homework:
\begin{enumerate}[itemsep=2pt]
  \item Turn each zero $a$ into the factor $(x - a)$ --- flip the sign.
  \item Raise each factor to its stated multiplicity.
  \item Check the degree: the multiplicities must add to the required degree.
  \item If a point is given, write $p(x) = a\,(\cdots)$, substitute the point,
        and solve for $a$.
\end{enumerate}
Worked live: degree 3, zero $3$ of multiplicity 2, zero $-5$, $p(0) = 90$
$\Rightarrow$ $45a = 90 \Rightarrow a = 2 \Rightarrow p(x) = 2(x-3)^2(x+5)$.
\end{skillbox}

\begin{skillbox}[Active Monitoring]{redbox}
While students work on guided notes and practice, circulate and check:
\begin{itemize}
  \item \textbf{Common error:} sign flip dropped --- reading the zero of
        $(x + 3)$ as $3$. Cold-call: ``What input makes that factor zero?''
  \item \textbf{Common error:} confusing multiplicity with the zero itself
        (``the zero is 2 with multiplicity $-1$''). Have them point at the
        exponent and then at the number inside the parentheses.
  \item \textbf{Common error:} answering with a number when a point is asked
        for, or vice versa. ``Zero or intercept? One is a number, one is a
        point.''
  \item \textbf{Common error:} in a difference table, subtracting the wrong
        direction (later minus earlier is the rule) or stopping at the second
        row out of habit from 2.1.
\end{itemize}
\end{skillbox}

\begin{skillbox}[Group Work \& Differentiation]{redbox}
Students work on the \textbf{Group Activity: Four Polynomials, Two Zeros} in
leveled tiers. All four formulas (A--D) share the zeros $-2$ and $1$ and
differ only in multiplicity, and all three graphs are pre-drawn --- no tier
draws anything.

\begin{itemize}
  \item \textbf{Tier R (Remediate):} Read the multiplicity of each zero off
        A--D, add them to get the degree, and label crossing vs.\ tangent at
        $x = 1$.
  \item \textbf{Tier A (Approaching Proficiency):} Match Graphs I, II, III to
        B, A, C (D is the leftover), justify one match from multiplicity
        rather than shape, and predict what D's graph would do
        (tangent at $-2$, crossing at $1$).
  \item \textbf{Tier E (Extension):} Construct $p$ with prescribed zeros,
        multiplicities, degree, and a point ($a = -2$,
        $p(x) = -2(x+3)^2(x-1)$); build a quartic tangent at two points and
        never crossing; and use a difference table to get $\deg q = 3$, then
        argue that a zero cannot have multiplicity 4.
\end{itemize}
\end{skillbox}

\begin{skillbox}[Individual Work \& Assessment]{redbox}
\textbf{Exit Ticket} (last 5--7 minutes, individual, collected):
\begin{enumerate}
  \item For $p(x) = (x-4)(x+1)^2$: table of zero, multiplicity, and
        crosses-or-tangent, plus the degree.
  \item Write a factored polynomial with a zero at $-3$ of multiplicity 1 and
        a zero at $2$ of multiplicity 2; state its degree.
  \item From a table of equally spaced values, compute first and second
        differences and name the degree.
\end{enumerate}
\textbf{Assessment note:} Item 1 is the diagnostic --- a student who writes
``tangent'' at $x = 4$ has memorized ``squared means tangent'' without
attaching it to the right factor. Item 2 catches the sign flip: a zero at
$-3$ must give $(x + 3)$. Item 3 is pure procedure; a miss there is an
arithmetic slip, not a concept gap.
\end{skillbox}

\begin{skillbox}[Reinforcement \& Extension]{goldbox}
\textbf{Homework:} 6 problems --- zeros, degree, and intercepts-as-points
from a factored cubic; two build-from-zeros items; a multiplicity table for a
degree-6 polynomial; a verify-the-zero computation feeding a
multiplicity read; a successive-difference table; and one AP-style
multiple-choice item on tangent vs.\ crossing.

\textbf{Extension (optional):} A degree-4 polynomial tangent at exactly two
points --- and an argument that three tangent points would force degree at
least 6, since each needs an even multiplicity.

\textbf{Preview:} Lesson 2.4, \textit{Polynomial Functions and Complex
Zeros}, takes up the loose thread: $p(x) = x^2 + 1$ has degree 2 and no
$x$-intercepts at all, so where did its zeros go?
\end{skillbox}

\begin{teachernote}[Warm-Up]
Five minutes, silent start, no calculators. Item 2 is the load-bearing one:
the zero-product property is the entire mechanism behind reading zeros off a
factored form, and if students cannot articulate ``one of the factors must be
zero'' they will not be able to justify anything in Part 2. Item 3 is not
just arithmetic --- it is the definition of a zero in disguise, so when you
review it, say the sentence ``so $x = 1$ \textit{is} a zero'' out loud and
leave it on the board for Part 1.
\end{teachernote}

\begin{teachernote}[Guided Notes]
Part 2 deserves the most clock. Write the equivalence as a double arrow and
physically trace it in both directions with your finger --- students who only
ever go factors-to-zeros collapse when asked to build a polynomial. The
sign-flip error in $(x+3)$ is the single most frequent mistake of this
lesson; catching it at the board costs thirty seconds and saves the homework.
In Part 3, resist stating the crossing/tangent rule first: run the sign
argument (an even power is never negative) and let the rule fall out, or it
becomes one more memorized pair of words. If time is short, compress Part 4
to a single board demonstration --- but do not skip it, it is exit-ticket
item 3 and it is the only place 1.5.A.6 appears in the unit.
\end{teachernote}

\begin{teachernote}[Group Activity]
The design point is that A--D all have the \textit{same} two zeros, so a
group cannot match graphs by reading intercepts --- multiplicity is the only
discriminator. Expect Tier A to try shape-matching first; push back with
``which factor is squared, and what does that force at that intercept?''
Groups can be moved between tiers mid-period without new materials since
every tier reads the same four formulas. Tier E item 3 is the bridge to
tomorrow: once they see that degree caps multiplicity, the question ``what if
the real multiplicities do not add up to the degree?'' becomes natural ---
bank it and open Lesson 2.4 with it.
\end{teachernote}

\begin{teachernote}[Exit Ticket]
Collected, completion credit only. Sort by item 1: a wrong crosses/tangent
pairing means Part 3 needs a five-minute revisit tomorrow before 2.4 adds
complex zeros on top. On item 2 accept either factor order but not
$(x - 3)(x + 2)^2$ --- that is the sign flip and it is worth naming to the
whole class. No notes; the whole ticket should take five minutes.
\end{teachernote}

\begin{teachernote}[Homework]
Problem 4 is the one to review first: part (a) is the definition of a zero
and part (b) hands them the factored form so the multiplicity read is
isolated from any factoring skill --- factoring technique is Lesson 2.6's
job, not tonight's. Problem 6 is the AP-style item; distractor C reverses the
multiplicities and distractor D makes both even, so both are worth naming
aloud. The extension's second half is a genuine proof-shaped argument; accept
any wording that counts even multiplicities to at least 6.
\end{teachernote}
\end{document}
